\chapter{PENDAHULUAN}
	\section{Latar Belakang Penelitian}
	\setlength{\leftskip}{0.7cm}
	Pemalsuan dokumen akademik seperti ijazah, transkip nilai, surat-menyurat dan sertifikasi kompetensi telah menjadi isu sangat krusial yang mengancam kredibilitas institusi pendidikan tinggi. Di Universitas Darussalam Gontor (UNIDA Gontor), pengelolaan dokumen akademik memerlukan standar kemanan tinggi untuk memastikan setiap dokumen tercatat dengan baik dan benar dan tidak dapat dimanipulasi. 

	Sistem verifikasi konvensional yang umum digunakan saat ini, termasuk di lingkungan UNIDA Gontor masih bergantung pada database terpusat dan proses manual,  yang memiliki kelemahan fundamental dalam hal keamanan data dan memerlukan waktu yang relatif lama. Database terpusat juga bertindak sebagai \textit{Single Point of Failure} (SPoF)\footcite{Awang2025}, yang mana jika server database terkena serangan siber atau mengalami kegagalan, maka seluruh proses akan terhenti, selain itu administrator juga memiliki akses penuh terhadap database sehingga memiliki kemampuan untuk mengubah atau memanipulasi data secara sepihak, sehingga integritas data tidak dapat dijamin sepenuhnya.

	Teknologi blockchain muncul sebagai paradigma baru yang menawarkan sifat \textit{immutability} (tidak dapat diubah) dan  transparansi melalui mekanisme penyimpanan data yang tersebar pada banyak node dalam jaringan \footcite{Zheng2017}. Namun, implementasi pada skala institusi pendidikan seringkali terbentur oleh hambatan ekonomi dan kompleksitas teknis. Pada blockchain publik, setiap interaksi data ke blockchain memerlukan biaya transaksi  atau disebut \textit{gas-fee} yang fluktuatif, yang menjadi beban finansial yang signifikan bagi universitas jika harus memverifikasi ribuan dokumen setiap tahunnya. Selain itu kurva pembelajaran yang terjal bagi para staf UNIDA Gontor seringkali menghambat untuk pengelolaan dompet kripto dan kunci kriptografi seringkali menghambat adopsi teknologi ini secara luas. 

	Dari sisi pengguna, tanda tangan digital pada umumnya seringkali direpresentasikan dalam bentuk deretan kode heksadesimal yang tidak dapat dibaca atau dipahami manusia (\textit{non-human readable}). Hal ini menimbulkan resiko keamanan di mana bagian arsip di UNIDA Gontor yang akan melakukan tanda tangan tidak mengetahui secara pasti data apa yang sedang mereka otorisasi tersebut. Implementasi protokol \textit{EIP-712} (\textit{Ethereum Typed Structured Data}) menjadi solusi krusial dalam penelitian ini karena memungkinkan penyajian data yang akan ditandatangani dalam format terstruktur yang mudah dibaca dan dipahami oleh manusia (\textit{human readable})\footcite{Bloemen2017}. Dengan protokol \textit{EIP-712}, biro dapat memverifikasi Nama, NIM dan Nomor Surat secara visual sebelum memberikan tanda tangan kriptografis, sehingga menjamin keakuratan otorisasi.  

	Guna mengatasi hambatan biaya transaksi, penelitian ini mengusulkan penggunaan \textit{Hyperledger Besu}\footcite{HyperledgerFoundation} sebagai infrastruktur \textit{private blockchain} dengan mekanisme \textit{zero-gas}, sehingga transaksi dapat divalidasi tanpa biaya (\textit{gas}) bagi pengirim. Selain itu, mekanisme verifikasi integritas dokumen mengadopsi dari penelitian Ramya dkk.\footcite{Ramya2026}, yaitu dengan membandingkan nilai \textit{hash} dokumen lokal dengan nilai \textit{hash} yang tercatat secara \textit{on-chain}. Dengan mengintegrasikan mekanisme tersebut dan manajemen kunci privat berdasarkan Pal dkk.\footcite{Pal2021}, melalui pendekatan \textit{local key management} (\textit{env/storage}), prototipe ini diharapkan mampu menyediakan solusi verifikasi dokumen yang aman, transparan, dan praktis untuk lingkungan akademik.



	Dalam perspektif keislaman, pengembangan sistem ini dapat dipandang sebagai implementasi dari prinsip maqashid syariah, khususnya dalam hal \textit{hifz al-‘aql} (menjaga akal). Konsep  \textit{hifz al-‘aql} tidak hanya sebatas berkaitan dengan perlindungan terhadap fungsi akal secara biologis, tetapi konsep ini juga mencangkup penjagaan keaslian, kebenaran, dan integritas ilmu pengetahuan yang dihasilkan oleh akal manusia. Tindakan pemalsuan dokumen akademik dapat dikategorikan sebagai bentuk penipuan (\textit{tadlis}) yang merusak nilai kebenaran  dan mencederai kehormatan ilmu itu sendiri. 

	Dengan demikian , prototipe ini hadir sebagai bentuk implementasi nilai-nilai kejujuran (\textit{As-Siddiq}) dan persaksian yang adil dalam ekosistem digital akademik. Sistem ini dirancang untuk memastikan bahwa setiap dokumen akademik yang diterbitkan dan diverifikasi memiliki keabsahan yang dapat dipertanggungjawabkan serta terhindar dari manipulasi data internal. Melalui penerapan mekanisme verifikasi yang aman, prototipe ini tidak hanya meningkatkan kepercayaan antar pihak, tetapi juga memperkuat transparansi dan akuntabilitas dalam pengelolaan dokumen akademik. Selain itu sistem ini diharapkan menjadi solusi yang relevan dalam menjawab tantangan digitalisasi, khususnya dalam menjaga integritas dan kredibilitas informasi akademik di UNIDA Gontor.

	\section{Rumusan Masalah}
	Berdasarkan permasalahan yang sudah di jelaskan di latar belakang, maka rumusan masalah dalam peneletian ini adalah:

	\begin{secenumerate}
		\item Sistem verifikasi dokumen akademik di UNIDA Gontor saat ini masih menggunakan database terpusat (\textit{centralized}), sehingga memiliki resiko \textit{Single Point of Failure} (SPoF) dan memungkinkan terjadinya manipulasi data internal. Selain itu, sistem yang ada belum menerapkan protokol penandatanganan data terstruktur seperti \textit{EIP-712} yang dapat meningkatkan keamanan sekaligus keterbacaan (\textit{human-readable}) tanda tangan digital. 

	\end{secenumerate}
	
	\section{Tujuan Penelitian}
	Berdasarkan rumusan masalah yang sudah diuraikan, maka tujuan penelitian ini sebagai berikut:

	\begin{secenumerate}
		\item Merancang sebuah prototipe \textit{API middleware} sistem verifikasi dokumen menggunakan standar \textit{EIP-712} yang  mampu membuktikan keaslian sebuah dokumen akademik UNIDA Gontor dan keabsahan tanda tangan pihak yang bersangkutan secara \textit{on-chain} di private blockchain. Serta, mengimplementasikan mekanisme transaksi gratis (\textit{gasless}) pada jaringan \textit{Hyperledger Besu} guna meniadakan biaya (\textit{gas}) saat transaksi berlangsung, sekaligus menerapkan skema manajemen kunci dengan memisahkan antara pemberi otoritas (\textit{signer}) dan juga pelaksana transaksi (\textit{relayer}) untuk menjamin keamanan dan integritas transaksi dokumen di lingkungan UNIDA Gontor.

	\end{secenumerate}

	\section{Batasan Masalah}
	Agar permasalahan yang sudah dirumuskan dapat terlaksana, maka penelitian ini memiliki batasan:

	\begin{secenumerate}
		\item Implementasi protokol fokus pada standar penandatanganan data terstruktur \textit{EIP-712}.
		\item Jaringan blockchain yang digunakan adalah \textit{Hyperledger Besu} dengan konsensus QBFT.
		\item Smart contract dikembangkan menggunakan bahasa \textit{Solidity 0.8.19} dan pustaka \textit{OpenZeppelin v4.9}.
		\item Sistem yang dikembangkan dalam penelitian ini masih berupa prototipe yang berfokus pada \textit{API Middleware} dan belum diimplementasikan secara langsung pada lingkungan operasional UNIDA Gontor, melainkan hanya digunakan sebagai model simulasi untuk menguji konsep dan fungsionalitas sistem.
		\item Penelitian ini tidak berfokus pada klasifikasi maupun pengelolaan seluruh jenis dokumen akademik di UNIDA Gontor, melainkan pada pengembangan mekanisme verifikasi berbasis parameter data yang bersifat generik sehingga dapat diadaptasikan pada berbagai jenis dokumen. 

	\end{secenumerate}
		
	\section{Manfaat Penelitian}
	Peneliti berharap adanya beberapa manfaat dari hasil penelitian ini, diantaranya sebagai berikut:
		\begin{secenumerate}
			\item Bagi Mahasiswa
			
			Penelitian ini memberikan referensi bagi mahasiswa dalam memahami dan mengimplementasikan protokol \textit{EIP-712} dalam jaringan private blockchain seperti \textit{Hyperledger besu} untuk menciptakan sistem verifikasi data yang aman dan transparan.

			\item Bagi Universitas
			
			Penelitian ini menjadi kontribusi akademik dan pengembangan riset terapan di bidang teknologi blockchain, serta meningkatkan kualitas penelitian yang berorientasi pada integritas data akademik.

			\item Bagi Masyarakat
			
			Penelitian ini menargetkan luaran prototipe sistem verifikasi bagi masyarakat dalam memvalidasi dokumen akademik.

			\item Bagi Peneliti
			
			Penelitian ini menjadi dasar dan referensi bagi peneliti selanjutnya dalam pengembangan manajemen identitas digital berbasis blockchain, serta penerapan  zero-gas pada jaringan privat untuk kebutuhan skala luas.
		\end{secenumerate}
	
	\section{Pembahasan Studi}
	Sistematika penulisan dalam penelitian ini terdiri dari lima bagian sebagai berikut:
	
	\noindent 
	\textbf{BAB 1: PENDAHULUAN}
	\begin{sysenumerate}
		\setcounter{sysenumeratei}{1}
		\item[] 
		\begin{sysenumerate}
			\item Latar Belakang Penelitian 
			\item Rumusan Masalah        
			\item Tujuan Penelitian 
			\item Batasan Masalah
			\item Manfaat Penelitian
			\item Pembahasan Studi
		\end{sysenumerate}
	\end{sysenumerate}
	
	\vspace{1em}

	\noindent 
	\textbf{BAB 2: TINJAUAN PUSTAKA}
	\begin{sysenumerate}
		\setcounter{sysenumeratei}{2}
		\item[] 
		\begin{sysenumerate}
			\item Penelitian Terdahulu         
			\item Landasan Teori
		\end{sysenumerate}
	\end{sysenumerate}
	
	\vspace{1em}
	
	\noindent 
	\textbf{BAB 3: METODOLOGI PENELITIAN}
	\begin{sysenumerate}
		\setcounter{sysenumeratei}{3} 
		\item[] 
		\begin{sysenumerate}
			\item Rencana Kegiatan 
			\item Alat dan Bahan Penelitian
			\item Tahapan Penelitian 
			\item Perancangan Sistem
			\item Sistematika Pembahasan
		\end{sysenumerate}
	\end{sysenumerate}
	
	\vspace{1em}
	
	\noindent 
	\textbf{BAB 4: HASIL DAN PEMBAHASAN}
	\begin{sysenumerate}
		\setcounter{sysenumeratei}{4}
		\item[] 
		\begin{sysenumerate}
			\item Hasil
			\item Pembahasan
		\end{sysenumerate}
	\end{sysenumerate}
		
	\vspace{1em}
		
	\noindent 
	\textbf{BAB 5: PENUTUP}
	\begin{sysenumerate}
	\setcounter{sysenumeratei}{5}
	\item[] 
		\begin{sysenumerate}
			\item Kesimpulan
			\item Saran
		\end{sysenumerate}
	\end{sysenumerate}